% ============================================================
%  AEGIS 综合设计报告 · 样式预置
%
%  设计取向：参照 IEEE Transactions / ACM TOPS 等期刊的排版惯例，
%  强调版面克制、图表自明、术语统一。
% ============================================================

% ---------- 页面与版心 ----------
\usepackage[a4paper,top=2.6cm,bottom=2.4cm,left=2.7cm,right=2.7cm,
            headheight=14pt,footskip=1.1cm]{geometry}

% ---------- 字体 ----------
\usepackage{amsmath,amssymb,amsthm}
\usepackage{mathtools}
\usepackage{bm}
% 正文西文使用 Times 系列，与中文宋体的字重相称
\setmainfont{Times New Roman}
\setsansfont{Arial}
\setmonofont{Consolas}[Scale=0.92]
% Times New Roman 不含 OMS 编码的数学符号字形（如 \vert 的可变尺寸形式），
% 显式将该符号族指向 Computer Modern，消除字体替换警告
\DeclareFontFamily{OMS}{TimesNewRoman(0)}{\skewchar\font=48}
\DeclareFontShape{OMS}{TimesNewRoman(0)}{m}{n}{<->ssub*cmsy/m/n}{}
\DeclareFontShape{OMS}{TimesNewRoman(0)}{b}{n}{<->ssub*cmsy/b/n}{}

% 中文字体：正文宋体、标题黑体、强调楷体
% 黑体本身即为粗笔画字体，无独立粗体字重，故显式指定 BoldFont 为自身，
% 避免 LaTeX 因找不到粗体变体而发出字体替换警告
\setCJKmainfont{SimSun}[BoldFont=SimHei, ItalicFont=KaiTi]
\setCJKsansfont{SimHei}[BoldFont=SimHei, ItalicFont=SimHei]
\setCJKmonofont{FangSong}[BoldFont=FangSong, ItalicFont=FangSong]

% ---------- 中文引号 ----------
% 正文中的引号统一使用全角弯引号字符（U+201C / U+201D），
% 由 CJK 字体渲染，以保证中文语境下的标点形态正确。
% 提供语义化命令以便统一调整。
\newcommand{\zh}[1]{“#1”}

% ---------- 颜色系统 ----------
\usepackage{xcolor}
\definecolor{aegisblue}{HTML}{1D4ED8}
\definecolor{aegisind}{HTML}{4F46E5}
\definecolor{aegisteal}{HTML}{0F766E}
\definecolor{aegisred}{HTML}{BE123C}
\definecolor{aegisamber}{HTML}{B45309}
\definecolor{aegisgray}{HTML}{475569}
\definecolor{rulegray}{HTML}{94A3B8}
\definecolor{panelbg}{HTML}{F8FAFC}
\definecolor{codebg}{HTML}{FAFBFD}
\definecolor{codekw}{HTML}{1D4ED8}
\definecolor{codestr}{HTML}{047857}
\definecolor{codecmt}{HTML}{64748B}

% ---------- 图表 ----------
\usepackage{graphicx}
\usepackage{booktabs}
\usepackage{multirow}
\usepackage{makecell}
\usepackage{array}
\usepackage{tabularx}
\usepackage{longtable}
\usepackage{threeparttable}
\usepackage{caption}
\usepackage{subcaption}
\usepackage{float}
\usepackage{rotating}

% 浮动体放置参数：放宽页面容纳浮动体的比例上限，
% 减少表格与图形被推迟到后续页面而在原处留下大片空白的情况
\renewcommand{\topfraction}{0.88}
\renewcommand{\bottomfraction}{0.72}
\renewcommand{\textfraction}{0.10}
\renewcommand{\floatpagefraction}{0.78}
\setcounter{topnumber}{3}
\setcounter{bottomnumber}{2}
\setcounter{totalnumber}{5}
% 浮动体与正文的间距
\setlength{\textfloatsep}{14pt plus 3pt minus 3pt}
\setlength{\intextsep}{12pt plus 3pt minus 3pt}
\setlength{\floatsep}{12pt plus 3pt minus 3pt}

% 表格行距略微放宽，提升可读性
\renewcommand{\arraystretch}{1.28}
\setlength{\tabcolsep}{6pt}

% 题注：小五号、加粗标签、居中
\captionsetup{
  font={small},
  labelfont={bf,color=aegisgray},
  labelsep=quad,
  skip=6pt,
  justification=centering
}
\captionsetup[table]{position=top,skip=5pt}
\captionsetup[figure]{position=bottom,skip=7pt}
\captionsetup[sub]{font={footnotesize},labelfont={bf}}

\renewcommand{\thetable}{\thesection--\arabic{table}}
\renewcommand{\thefigure}{\thesection--\arabic{figure}}
% 图表编号按章重新计数。ctexart 的 section 是顶层章节单位，
% 但 LaTeX 默认不会在 section 处重置 table/figure 计数器，
% 需显式声明依赖关系，否则编号会跨章连续累加（如出现"表 6--17"）
\counterwithin{table}{section}
\counterwithin{figure}{section}
\counterwithin{equation}{section}

% ---------- TikZ / PGFPlots ----------
\usepackage{tikz}
\usepackage{pgfplots}
\pgfplotsset{compat=1.18}
\usetikzlibrary{
  arrows.meta, positioning, calc, fit, backgrounds, shapes.geometric,
  shapes.misc, shapes.symbols, decorations.pathreplacing,
  decorations.pathmorphing, patterns, matrix, chains, shadows.blur, trees
}
\usepgfplotslibrary{groupplots, fillbetween, statistics}
% 甘特图里程碑等符号
\usepackage{pifont}

% 统一的绘图基调：细网格、无外框、克制的配色
\pgfplotsset{
  every axis/.append style={
    line width=0.55pt,
    tick style={line width=0.5pt, color=rulegray},
    label style={font=\small},
    tick label style={font=\footnotesize},
    title style={font=\small\bfseries},
    legend style={
      font=\footnotesize,
      draw=rulegray!60,
      line width=0.4pt,
      fill=white,
      fill opacity=0.92,
      text opacity=1,
      rounded corners=1.5pt
    },
    grid style={rulegray!30, line width=0.35pt, dashed},
    axis line style={rulegray!85}
  },
  % 期刊图常用：仅保留左侧与底部坐标轴
  scientific/.style={
    axis lines=left,
    ymajorgrids=true,
    enlarge x limits=0.06,
    enlarge y limits={upper, value=0.12}
  }
}

% TikZ 通用节点样式
\tikzset{
  >=Latex,
  block/.style={
    draw=aegisblue!70, fill=aegisblue!6, line width=0.6pt,
    rounded corners=2.2pt, align=center, font=\small,
    inner sep=5pt, minimum height=8mm
  },
  blockalt/.style={
    draw=aegisteal!70, fill=aegisteal!6, line width=0.6pt,
    rounded corners=2.2pt, align=center, font=\small,
    inner sep=5pt, minimum height=8mm
  },
  blockred/.style={
    draw=aegisred!75, fill=aegisred!7, line width=0.7pt,
    rounded corners=2.2pt, align=center, font=\small,
    inner sep=5pt, minimum height=8mm
  },
  astnode/.style={
    draw=aegisblue!65, fill=aegisblue!8, line width=0.55pt,
    rounded corners=2pt, align=center, font=\scriptsize,
    inner sep=3pt, minimum width=13mm, minimum height=5.4mm
  },
  astinj/.style={
    draw=aegisred!80, fill=aegisred!10, line width=0.85pt,
    rounded corners=2pt, align=center, font=\scriptsize\bfseries,
    inner sep=3pt, minimum width=13mm, minimum height=5.4mm
  },
  lbl/.style={font=\scriptsize, color=aegisgray},
  flow/.style={->, line width=0.6pt, color=aegisgray!85},
  flowred/.style={->, line width=0.9pt, color=aegisred!85}
}

% ---------- 代码环境 ----------
\usepackage{listings}
\lstset{
  basicstyle=\ttfamily\scriptsize,
  keywordstyle=\color{codekw}\bfseries,
  stringstyle=\color{codestr},
  commentstyle=\color{codecmt}\itshape,
  numberstyle=\ttfamily\tiny\color{rulegray},
  numbers=left,
  numbersep=7pt,
  backgroundcolor=\color{codebg},
  frame=single,
  framerule=0.4pt,
  rulecolor=\color{rulegray!45},
  breaklines=true,
  breakatwhitespace=false,
  showstringspaces=false,
  tabsize=2,
  xleftmargin=14pt,
  xrightmargin=2pt,
  aboveskip=8pt,
  belowskip=6pt,
  captionpos=b,
  escapeinside={(*@}{@*)}
}

\lstdefinestyle{javastyle}{language=Java}
\lstdefinestyle{sqlstyle}{
  language=SQL,
  morekeywords={UNION,SLEEP,BENCHMARK,LOAD_FILE,EXTRACTVALUE,CAST,VARCHAR},
  emph={SLEEP,BENCHMARK,LOAD_FILE,EXTRACTVALUE,UNION},
  emphstyle=\color{aegisred}\bfseries
}

% 代码清单题注汉化并按章编号。
% lstlisting 计数器由 listings 宏包在文档开始时创建，
% 故延迟至 \AtBeginDocument 阶段配置
\renewcommand{\lstlistingname}{代码}
\renewcommand{\lstlistlistingname}{代码目录}
\AtBeginDocument{%
  \counterwithin{lstlisting}{section}%
  \renewcommand{\thelstlisting}{\thesection--\arabic{lstlisting}}%
}

% ---------- 算法 ----------
\usepackage[ruled,linesnumbered,algosection]{algorithm2e}
\SetAlCapFnt{\small\bfseries\color{aegisgray}}
\SetAlCapNameFnt{\small}
\SetAlgoCaptionSeparator{\quad}
\SetKwInOut{Input}{输入}
\SetKwInOut{Output}{输出}
\SetKwFor{For}{for}{do}{end}
\SetKwFor{While}{while}{do}{end}
\SetKwIF{If}{ElseIf}{Else}{if}{then}{else if}{else}{end}
\SetKwRepeat{Repeat}{repeat}{until}

% ---------- 定理环境 ----------
\theoremstyle{definition}
\newtheorem{definition}{定义}[section]
\newtheorem{property}{性质}[section]
\theoremstyle{plain}
\newtheorem{proposition}{命题}[section]
\newtheorem{theorem}{定理}[section]
\theoremstyle{remark}
\newtheorem{remark}{注}[section]

% ---------- 版式细节 ----------
\usepackage{fancyhdr}
\usepackage{titlesec}
\usepackage{enumitem}
\usepackage{footmisc}
\usepackage{tcolorbox}
\tcbuselibrary{skins,breakable}

\setlist{itemsep=1pt, parsep=1pt, topsep=3pt, leftmargin=*}

% 章节标题：黑体、适度留白
\ctexset{
  section={
    format={\Large\sffamily\bfseries},
    beforeskip={16pt plus 2pt minus 2pt},
    afterskip={10pt plus 1pt}
  },
  subsection={
    format={\large\sffamily\bfseries},
    beforeskip={12pt plus 2pt minus 1pt},
    afterskip={7pt plus 1pt}
  },
  subsubsection={
    format={\normalsize\sffamily\bfseries},
    beforeskip={10pt plus 1pt},
    afterskip={5pt}
  }
}

% 页眉页脚
\pagestyle{fancy}
\fancyhf{}
\renewcommand{\headrulewidth}{0.4pt}
\renewcommand{\headrule}{\hbox to\headwidth{\color{rulegray!70}\leaders\hrule height \headrulewidth\hfill}}
\fancyhead[L]{\footnotesize\color{aegisgray}\sffamily AEGIS：基于抽象语法树结构指纹的 Web 攻击语义检测}
\fancyhead[R]{\footnotesize\color{aegisgray}\thepage}

\fancypagestyle{plain}{
  \fancyhf{}
  \renewcommand{\headrulewidth}{0pt}
  \fancyfoot[C]{\footnotesize\color{aegisgray}\thepage}
}

% 要点提示框
\newtcolorbox{keybox}[1][]{
  enhanced, breakable,
  colback=panelbg, colframe=aegisblue!55,
  boxrule=0.5pt, left=8pt, right=8pt, top=6pt, bottom=6pt,
  arc=2pt,
  borderline west={2.2pt}{0pt}{aegisblue},
  fontupper=\small,
  #1
}

\newtcolorbox{warnbox}[1][]{
  enhanced, breakable,
  colback=aegisred!3, colframe=aegisred!45,
  boxrule=0.5pt, left=8pt, right=8pt, top=6pt, bottom=6pt,
  arc=2pt,
  borderline west={2.2pt}{0pt}{aegisred},
  fontupper=\small,
  #1
}

% ---------- 超链接（置于最后加载）----------
\usepackage[
  colorlinks=true,
  linkcolor=aegisblue,
  citecolor=aegisteal,
  urlcolor=aegisind,
  bookmarksnumbered=true,
  pdfstartview=FitH
]{hyperref}
\usepackage{cleveref}
\crefname{figure}{图}{图}
\crefname{table}{表}{表}
\crefname{section}{第}{第}
\crefname{equation}{式}{式}
\crefname{algorithm}{算法}{算法}
\crefname{definition}{定义}{定义}
\crefname{property}{性质}{性质}

% ---------- 常用宏 ----------
\newcommand{\aegis}{\textsc{Aegis}}
\newcommand{\code}[1]{\texttt{\small #1}}
\newcommand{\term}[1]{\textbf{#1}}
% 数学记号
\newcommand{\fp}{\ensuremath{\mathcal{F}}}          % 指纹函数
\newcommand{\Norm}{\ensuremath{\mathcal{N}}}        % 归一化算子
\newcommand{\Ser}{\ensuremath{\mathcal{S}}}         % 序列化算子
\newcommand{\Tree}{\ensuremath{\mathcal{T}}}        % 语法树
\newcommand{\Sig}{\ensuremath{\Sigma}}              % 签名多重集
\newcommand{\jac}{\ensuremath{J}}                   % Jaccard 相似度
